\documentclass{article}
\usepackage{mathtools} 
\usepackage{fontspec}
\usepackage[UTF8]{ctex}
\usepackage{amsthm}
\usepackage{mdframed}
\usepackage{xcolor}
\usepackage{amssymb}
\usepackage{amsmath}


% 定义新的带灰色背景的说明环境 zremark
\newmdtheoremenv[
  backgroundcolor=gray!10,
  % 边框与背景一致，边框线会消失
  linecolor=gray!10
]{zremark}{说明}

% 通用矩阵命令: \flexmatrix{矩阵名}{元素符号}{行数}{列数}
\newcommand{\flexmatrix}[4]{
  \[
  #1 = \begin{pmatrix}
    #2_{11}     & #2_{12}     & \cdots & #2_{1#4}   \\
    #2_{21}     & #2_{22}     & \cdots & #2_{2#4}   \\
    \vdots      & \vdots      & \ddots & \vdots     \\
    #2_{#31}    & #2_{#32}    & \cdots & #2_{#3#4}
  \end{pmatrix}
  \]
}

% 简化版命令（默认矩阵名为A，元素符号为a）: \quickmatrix{行数}{列数}
\newcommand{\quickmatrix}[2]{\flexmatrix{A}{a}{#1}{#2}}

\begin{document}
\title{习题 6.2}
\author{张志聪}
\maketitle

\section*{9}

若$f,g$满足：
\begin{itemize}
  \item (i) $\lim\limits_{x \to +\infty} f(x) = 0, \lim\limits_{x \to +\infty} g(x) = 0$。
  \item (ii) 在$U(+\infty)$上，$f, g$可导，且$g'(x) \neq 0$。
  \item (iii) $\lim\limits_{x \to +\infty} \frac{f'(x)}{g'(x)} = A$，则
        \begin{align*}
          \lim\limits_{x \to +\infty} \frac{f(x)}{g(x)} = \lim\limits_{x \to +\infty} \frac{f'(x)}{g'(x)} = A
        \end{align*}
\end{itemize}

\textbf{证明：}

变量替换：令$t = \frac{1}{x}$，
于是$t \to 0^+$时，$x \to +\infty$，且
\begin{itemize}
  \item (i) $\lim\limits_{t \to 0^+} f(\frac{1}{t}) = 0, \lim\limits_{t \to 0^+} g(\frac{1}{t}) = 0$。
  \item (ii) 在$t \in U_+^o(0)$上，$f(\frac{1}{t}), g(\frac{1}{t})$可导，
        且$(g(\frac{1}{t}))' = - \frac{g'(\frac{1}{t})}{t^2} \neq 0$。
  \item (iii) $\lim\limits_{t \to 0^+} \frac{(f(\frac{1}{t}))'}{(g(\frac{1}{t}))'} = \lim\limits_{t \to 0^+} \frac{f'(\frac{1}{t})}{g'(\frac{1}{t})} = A$。
\end{itemize}

利用定理6.7可知
\begin{align*}
  \lim\limits_{x \to +\infty} \frac{f(x)}{g(x)}
   & = \lim\limits_{t \to 0^+} \frac{f(\frac{1}{t})}{g(\frac{1}{t})}       \\
   & = \lim\limits_{t \to 0^+} \frac{(f(\frac{1}{t}))'}{(g(\frac{1}{t}))'} \\
   & = A
\end{align*}


\section*{10}

提示：分段
\end{document}